% =====================================================================
% Module 3: Xử lý dữ liệu (Processing Layer)
% ---------------------------------------------------------------------
% Cách sử dụng:
%   1. Đặt file .tex này cùng thư mục với báo cáo chính, hoặc \input{} nó.
%   2. Mở file "module3_diagrams.html" bằng trình duyệt, chụp màn hình từng
%      biểu đồ và lưu thành ảnh (PNG/PDF) đặt cùng thư mục với tên:
%      module3_hinh3_1.png  (Luồng xử lý Silver theo từng ngày)
%      module3_hinh3_2.png  (So sánh bộ cleaner của 3 trang web)
%      module3_hinh3_3.png  (Biến đổi sơ đồ hình sao Gold)
%   3. Gỡ bỏ khối "lstlisting setup" bên dưới nếu báo cáo đã có package
%      "listings" và định nghĩa kiểu Python.
% =====================================================================

% ---- Vietnamese encoding (xoa neu da co trong preamble) -------------
\usepackage[utf8]{inputenc}
\usepackage[T5]{fontenc}
% ---- lstlisting setup (xoa neu da co trong preamble) ----------------
\usepackage{listings}
\usepackage{xcolor}
\lstdefinestyle{pystyle}{
    language=Python,
    basicstyle=\ttfamily\footnotesize,
    keywordstyle=\color{blue}\bfseries,
    stringstyle=\color{teal},
    commentstyle=\color{gray}\itshape,
    showstringspaces=false,
    breaklines=true,
    frame=single,
    rulecolor=\color{lightgray},
    backgroundcolor=\color{white},
    tabsize=2,
    columns=fullflexible,
    keepspaces=true,
}
\lstdefinestyle{sqlstyle}{
    language=SQL,
    basicstyle=\ttfamily\footnotesize,
    keywordstyle=\color{blue}\bfseries,
    stringstyle=\color{teal},
    commentstyle=\color{gray}\itshape,
    showstringspaces=false,
    breaklines=true,
    frame=single,
    rulecolor=\color{lightgray},
    backgroundcolor=\color{white},
    tabsize=2,
    columns=fullflexible,
    keepspaces=true,
}
\lstset{style=pystyle}
% ---------------------------------------------------------------------

\usepackage{graphicx}
\usepackage{booktabs}
\usepackage{longtable}
\usepackage{hyperref}


\section{Module 3: Xử lý dữ liệu (Processing Layer)}
\label{sec:processing}

Processing Layer là nơi thực hiện các phép biến đổi dữ liệu từ
Bronze (thô, \texttt{.jsonl.gz}, cấu trúc không đồng nhất) sang Silver
(\texttt{.parquet}, cấu trúc chuẩn \texttt{SilverJobItem}), rồi tiếp tục
biến đổi sang Gold (sơ đồ hình sao - star schema trên MotherDuck) và phục vụ (serving: Supabase,
Qdrant). Khác với Module~2 tập trung vào \emph{vị trí} lưu trữ (S3 buckets,
phân vùng - partitioning), Module~3 tập trung vào \emph{logic} biến đổi.

Các thành phần chính của Processing Layer:
\begin{itemize}
  \item \textbf{Silver cleaning pipeline} — khung điều phối (orchestrator) đọc
    Bronze theo từng ngày (day-by-day), chạy dãy hàm làm sạch (cleaner) theo thứ tự, ép schema, tải lên Silver.
  \item \textbf{Các hàm làm sạch dùng chung (Common cleaners)} — hơn 15 hàm dùng chung giữa 3 trang web
    (TopCV, ITViec, VietnamWorks), mỗi hàm chịu trách nhiệm làm sạch một trường dữ liệu (field).
  \item \textbf{Seed taxonomy} — 13 tệp CSV (ưu tiên Google Sheets, tệp cục bộ làm dự phòng - local fallback) dùng để ánh xạ (map) giá trị thô sang nhãn chuẩn hóa (canonical label).
  \item \textbf{Bộ trích xuất lai FlashText + regex (FlashText + regex hybrid extractor)} — trích xuất kỹ năng (skill) từ văn bản với tốc độ nhanh và độ chính xác cao.
  \item \textbf{Bronze validation} — kiểm tra chất lượng dữ liệu bằng Great Expectations trước khi đưa vào Silver.
  \item \textbf{Gold SQL transformations} — 7 câu lệnh SQL xây dựng sơ đồ hình sao (loại bỏ trùng lặp trung gian - staging dedup $\to$ bảng sự kiện - fact + chiều thời gian - dim\_date + các bảng con - child tables).
  \item \textbf{Bảng điều khiển (Business dashboards)} — báo cáo dạng HTML tổng hợp dữ liệu từ Silver/Bronze.
\end{itemize}

Tổng quan luồng xử lý được thể hiện trong Hình~\ref{fig:silver-pipeline}.

\begin{figure}[htbp]
  \centering
  % Screenshot: module3_diagrams.html diagram 1 (pipeline flow)
  \includegraphics[width=\textwidth]{module3_hinh3_1.png}
  \caption{Luồng xử lý Silver theo từng ngày (day-by-day): đọc Bronze mỗi ngày, lọc dòng
  thừa, chạy dãy hàm làm sạch (cleaner), ép schema, tải lên Silver S3.}
  \label{fig:silver-pipeline}
\end{figure}

% ============================================================
\subsection{Khung Silver cleaning pipeline}
\label{subsec:silver-pipeline-framework}

% ============================================================
\subsubsection{Bộ điều phối (Orchestrator): \texttt{run\_pipeline()}}

Hàm \texttt{run\_pipeline()} trong \texttt{common/pipeline.py} là điểm
bắt đầu (entry point) duy nhất cho tất cả 3 trang web (site). Nó đọc dữ liệu Bronze theo từng ngày riêng biệt
(từ \texttt{from\_date} đến \texttt{to\_date}), áp dụng các hàm làm sạch (cleaning), ép schema,
rồi tải lên Silver S3.

\begin{lstlisting}[style=pystyle, caption={Vòng lặp điều phối đọc theo ngày}, label={lst:run-pipeline}]
def run_pipeline(site, clean_fn, args):
    """Read Bronze day-by-day -> clean -> upload to Silver S3."""
    start_dt = datetime.strptime(args.from_date, "%Y-%m-%d")
    end_dt = datetime.strptime(args.to_date, "%Y-%m-%d")
    current_dt = start_dt
    while current_dt <= end_dt:
        date_str = current_dt.strftime("%Y-%m-%d")
        lazy = get_jobs_data_from_bronze(
            site=site, entity_name=args.entity_name,
            from_date=date_str, to_date=date_str,
        )
        if lazy is None:
            current_dt += timedelta(days=1)
            continue
        df = lazy.collect()
        df = filter_essential_rows(df)
        cleaned = clean_fn(df)                     # site-specific
        cleaned = enforce_silver_schema(cleaned, SILVER_SCHEMA)
        if args.export_parquet:
            # Debug dump to local parquet for visual inspection
            cleaned.write_parquet(SILVER_DEBUG_DIR / f"{site}_cleaned_{date_str}.parquet")
        if not args.no_save:
            s3_key = upload_silver_parquet(
                df=cleaned, entity_name=args.entity_name,
                source_site=site, date_str=date_str,
            )
        current_dt += timedelta(days=1)
\end{lstlisting}

Thiết kế đọc theo ngày (không đọc hết toàn bộ khoảng thời gian cùng một lúc) có mục đích:
chạy lại (re-run) chỉ cần xử lý các ngày cần thiết, và nếu một ngày bị lỗi (fail),
các ngày khác vẫn hoạt động bình thường do mỗi ngày là độc lập.

% ============================================================
\subsubsection{Lọc dòng thừa: \texttt{filter\_essential\_rows()}}

Trước khi chạy các bước làm sạch (để không lãng phí tài nguyên CPU cho các dòng rác),
đường ống (pipeline) loại bỏ các dòng thiếu các trường dữ liệu (field) thiết yếu:

\begin{lstlisting}[style=pystyle, caption={Lọc dòng thiếu field thiết yếu}, label={lst:filter-essential}]
ESSENTIAL_COLUMNS = ("job_title", "company_name")

def filter_essential_rows(df):
    """Drop rows where any essential column is null."""
    before = df.height
    cols_present = [c for c in ESSENTIAL_COLUMNS if c in df.columns]
    if cols_present:
        df = df.filter(~pl.any_horizontal(pl.col(c).is_null() for c in cols_present))
    after = df.height
    if before - after:
        logger.info("Dropped %d/%d rows missing %s", before-after, before, cols_present)
    return df
\end{lstlisting}

% ============================================================
\subsubsection{Ép kiểu cấu trúc (Schema): \texttt{enforce\_silver\_schema()}}

Sau khi làm sạch (cleaning), DataFrame có thể thiếu cột (do trang web không thu thập được trường dữ liệu đó)
hoặc sai kiểu dữ liệu (ví dụ Chuỗi thay vì Số thực). Hàm này đưa DataFrame về đúng
cấu trúc chuẩn của \texttt{SilverJobItem}:

\begin{lstlisting}[style=pystyle, caption={Chuyển đổi DataFrame về cấu trúc chuẩn SilverJobItem}, label={lst:enforce-schema}]
def enforce_silver_schema(df, schema):
    """Cast existing, add missing with null, drop extra, reorder."""
    existing = set(df.columns)
    exprs = []
    for col, dtype in schema.items():
        if col in existing:
            src_dtype = df.schema[col]
            # Polars does not support String -> Boolean; bridge via Int64
            if dtype == pl.Boolean and src_dtype == pl.String():
                exprs.append(pl.col(col).cast(pl.Int64, strict=False).cast(pl.Boolean))
            else:
                exprs.append(pl.col(col).cast(dtype, strict=False))
        else:
            exprs.append(pl.lit(None).cast(dtype).alias(col))
    return df.select(exprs)  # return in schema order
\end{lstlisting}

Schema \texttt{SILVER\_SCHEMA} được sinh tự động từ lớp dữ liệu (dataclass)
\texttt{SilverJobItem} thông qua hàm \texttt{silver\_schema\_to\_polars()} —
khi thêm hoặc bớt trường dữ liệu (field), chỉ cần sửa đổi dataclass mà không cần chỉnh sửa chuỗi schema bằng tay.

% ============================================================
\subsection{Các hàm làm sạch dùng chung (Common cleaners)}
\label{subsec:common-cleaners}

Mỗi hàm làm sạch dùng chung (common cleaner) nhận đầu vào là \texttt{pl.DataFrame}, thực hiện xử lý và trả về \texttt{pl.DataFrame}
được bổ sung thêm cột dữ liệu đã làm sạch (cleaned). Hàm đặc thù của từng trang web \texttt{clean\_<site>\_jobs()} sẽ kết hợp
chúng theo một thứ tự cố định. Hình~\ref{fig:site-compare} so sánh bộ hàm làm sạch của 3 trang web.

\begin{figure}[htbp]
  \centering
  % Screenshot: module3_diagrams.html diagram 2 (site comparison)
  \includegraphics[width=\textwidth]{module3_hinh3_2.png}
  \caption{So sánh bộ làm sạch của 3 trang web: shared = các hàm dùng chung (common cleaner),
  site-specific = hàm bọc (wrapper) kèm theo extra\_noise\_patterns riêng của từng trang,
  skip = bỏ qua do cột dữ liệu không tồn tại.}
  \label{fig:site-compare}
\end{figure}

% ============================================================
\subsubsection{Chuẩn hóa URL: \texttt{clean\_job\_url()}}

Tạo trường \texttt{unique\_url} làm khóa loại bỏ trùng lặp (deduplication) giữa các trang web.
Đường dẫn URL của ITViec có phần đuôi số dạng numeric (\texttt{-12345.html}) nên cần một biểu thức chính quy (regex) loại bỏ đuôi số riêng:

\begin{lstlisting}[style=pystyle, caption={URL normalization + unique\_url dedup key}, label={lst:clean-url}]
def clean_job_url(df, column_name="job_url",
                  unique_column_name="unique_url",
                  source_site_column="source_site"):
    cleaned_url = clean_url_expr(pl.col(column_name))
    source_site = _source_site_expr(df, source_site_column)
    return df.with_columns(
        cleaned_url.alias(column_name),
        build_unique_url_expr(cleaned_url, source_site).alias(unique_column_name),
    )
\end{lstlisting}

Hàm \texttt{build\_unique\_url\_expr()} sẽ ghép \texttt{source\_site} +
\texttt{job\_url} đã được loại bỏ các phần tham số truy vấn (query) và phân mảnh (fragment). Đối với ITViec, hàm áp dụng thêm
\texttt{regexp\_replace(url, '-[0-9]+/?\$', '')} để loại bỏ phần đuôi số.

% ============================================================
\subsubsection{Chuỗi xử lý tên công ty: \texttt{main\_clean\_company()}}

Chuỗi xử lý gồm 5 bước, chạy theo \textbf{thứ tự bắt buộc} vì \texttt{clean\_2} phải
chạy trước \texttt{clean\_1} (nếu không, phần văn bản chứa trong dấu ngoặc đơn sẽ bị sót lại):

\begin{lstlisting}[style=pystyle, caption={Chuỗi làm sạch tên công ty gồm 5 bước}, label={lst:clean-company}]
def main_clean_company(df, column_name="company_name"):
    """NFC -> clean_2 (strip URLs, brackets) -> clean_1 (strip legal forms)
    -> map_canonical_company (join seed) -> coalesce -> drop intermediate."""
    df = normalize_unicode(df, column_name)          # NFC normalize
    df = clean_2(df, column_name)                    # strip URLs, "(SHB)" brackets
    df = clean_1(df, column_name, new_column_name=column_name)  # strip "Co., Ltd"
    df = map_canonical_company(df, source_col=column_name)      # join company_mapping.csv
    df = df.with_columns(
        pl.coalesce(pl.col("company_name_canonical"), pl.col(column_name))
        .alias("clean_company_name")                  # unmapped = fallback to raw
    )
    df = df.drop([column_name, "company_name_canonical"])
    return df
\end{lstlisting}

Hàm \texttt{map\_canonical\_company()} liên kết (join) với tệp \texttt{company\_mapping.csv}
để ánh xạ các tên biến thể (ví dụ ``FPT Telecom'', ``FPT'') sang tên chính thức chuẩn hóa
(``FPT Corporation''). Khi không khớp (match), \texttt{company\_name\_canonical = NULL}
để các bước xử lý phía sau (downstream) phân biệt được với chuỗi trống.

% ============================================================
\subsubsection{Đường ống làm sạch chức danh: \texttt{process\_job\_title\_pipeline()}}

Quy trình gồm 4 bước: (1) loại bỏ dấu tiếng Việt, (2) trích xuất từ khóa cấp bậc + vai trò (role),
(3) làm sạch chức danh (title) gốc, (4) xóa các cột dữ liệu trung gian:

\begin{lstlisting}[style=pystyle, caption={Job title: trích keyword + clean}, label={lst:title-pipeline}]
def process_job_title_pipeline(df, title_col="job_title"):
    df = remove_accents_col(df, column_name=title_col,
                            new_column_name="job_title_no_accent")
    df = extract_job_keywords(df, title_no_accent_col="job_title_no_accent",
                              keyword_col="job_title_special_keywords")
    df = clean_job_title(df, column_name=title_col,
                         new_column_name="clean_job_title")
    df = df.drop("job_title_no_accent")
    return df
\end{lstlisting}

Hàm \texttt{extract\_job\_keywords()} thực hiện trích xuất từ khóa thành 3 nhóm: \textbf{Level} (cấp bậc: Intern,
Fresher, Junior, Senior, Manager/Lead), \textbf{Domain} (lĩnh vực chuyên môn: Data, Business,
ML, AI, Database, Analytics), \textbf{Action} (vai trò công việc: Scientist, Engineer,
Analyst, Administrator). Mỗi từ khóa được kiểm tra độc lập để nhận diện được các chức danh kết hợp nhiều vai trò (multi-role)
(ví dụ ``Data Analytics and Engineer'' $\to$ [``Data'', ``Analytics'',
``Engineer'']).

% ============================================================
\subsubsection{Xử lý phân tích mức lương: \texttt{clean\_salary()}}

Hàm có cấu trúc logic phức tạp nhất — thực hiện phân tích cú pháp chuỗi mô tả mức lương tự do thành 2 trường số riêng biệt
\texttt{min\_monthly\_salary} / \texttt{max\_monthly\_salary} (VND/tháng):

\begin{lstlisting}[style=pystyle, caption={Xử lý phân tích lương: đơn vị + thời gian + khoảng lương}, label={lst:clean-salary}]
def clean_salary(df, column_name="salary",
                 min_col="min_monthly_salary", max_col="max_monthly_salary"):
    """Parse salary into min/max (VND/month).
    Currency: USD -> *25000; k -> *1000; tr/trieu -> *1M; ty -> *1B.
    Time: day -> *30; week -> *4; year -> *0.083.
    Range: '10 - 20' -> min=10, max=20; 'over X' -> max=null; 'upto X' -> min=null.
    Single number <1000 without unit -> implicit million VND."""
    raw = pl.col(column_name).cast(pl.String)
    salary = raw.str.to_lowercase()
    is_usd = salary.str.contains(r"(usd|\$)")
    has_million = salary.str.contains(r"(tr|trieu|m)")
    has_thousand = salary.str.contains(r"(nghin|k)")
    has_day = salary.str.contains("ngay"); has_week = salary.str.contains("tuan")
    # ... (classification flags for each case)
    unit_mult = (
        pl.when(has_thousand).then(pl.lit(1_000.0))
        .when(has_million).then(pl.lit(1_000_000.0))
        .when(is_usd).then(pl.lit(float(USD_TO_VND_RATE)))
        .otherwise(pl.lit(1.0))
    )
    time_mult = (
        pl.when(has_day).then(pl.lit(30.0))
        .when(has_week).then(pl.lit(4.0))
        .otherwise(pl.lit(1.0))
    )
    # Extract first and second number, multiply by unit_mult * time_mult
    # ...
    return df.with_columns([min_expr.alias(min_col), max_expr.alias(max_col)])
\end{lstlisting}

% ============================================================
\subsubsection{Làm sạch ngành nghề: \texttt{apply\_industry\_cleaning()}}

Quy trình (pattern): tách chuỗi (split) $\to$ phân rã dòng (explode) $\to$ chuẩn hóa (normalize) $\to$ ánh xạ (map) $\to$ nhóm lại (group back).
TopCV sử dụng dấu phẩy phân cách (\texttt{sep=","}); trong khi ITViec/VietnamWorks chứa chuỗi đơn (\texttt{sep=None}):

\begin{lstlisting}[style=pystyle, caption={Ngành nghề: tách, ánh xạ taxonomy, và nhóm lại}, label={lst:clean-industry}]
def apply_industry_cleaning(df, taxonomy_df, sep=None):
    df = df.with_row_index("temp_row_idx")
    if sep:
        df_exploded = df.with_columns(
            pl.col("job_industry").cast(pl.String).str.split(sep).alias("industry_split"))
    else:
        df_exploded = df.with_columns(
            pl.col("job_industry").cast(pl.String)
            .map_elements(lambda x: [x] if x else [], return_dtype=pl.List(pl.String))
            .alias("industry_split"))
    df_exploded = df_exploded.explode("industry_split").with_columns(
        pl.col("industry_split").str.strip_chars()
    ).filter(pl.col("industry_split").is_not_null() & (pl.col("industry_split") != ""))
    # Normalize: strip accents, lowercase -> map taxonomy -> group back
    df_mapped = df_exploded.with_columns(
        pl.col("industry_split").map_elements(remove_vietnamese_accents, return_dtype=pl.String)
        .str.to_lowercase().alias("industry_split_norm"))
    mapping_expr = build_industry_mapping_expr(taxonomy_df, "industry_split_norm")
    df_mapped = df_mapped.with_columns(mapping_expr.alias("mapped_industry"))
    # Unmapped -> "Others"; track originals in job_industry_unmapped
    df_handled = df_mapped.with_columns(
        pl.when(pl.col("mapped_industry").is_null())
          .then(pl.col("industry_split")).otherwise(pl.lit(None)).alias("unmapped_element"),
        pl.col("mapped_industry").fill_null("Others").alias("clean_industry_element"),
    )
    df_cleaned = df_handled.group_by("temp_row_idx").agg(
        pl.col("clean_industry_element").drop_nulls().unique().alias("job_industry_clean"),
        pl.col("unmapped_element").drop_nulls().unique().alias("job_industry_unmapped"))
    return df.join(df_cleaned, on="temp_row_idx", how="left").drop("temp_row_idx")
\end{lstlisting}

Cột \texttt{job\_industry\_unmapped} lưu giữ các giá trị không khớp với
bộ phân loại (taxonomy) để nhóm phát triển dữ liệu đánh giá và cập nhật tệp \texttt{industry\_taxonomy.csv}.

% ============================================================
\subsubsection{Phúc lợi (Benefit) và Yêu cầu (Requirement): Trích xuất FlashText}

Cả hai trường dữ liệu phúc lợi và yêu cầu tuyển dụng đều áp dụng chung một quy trình: dùng hàm \texttt{clean\_text()}
$\to$ loại bỏ dấu $\to$ viết thường $\to$ so khớp từ khóa bằng FlashText (FlashText keyword matching).
Điểm khác biệt ở cột yêu cầu (Requirement) là nó áp dụng 8 taxonomy riêng biệt:

\begin{lstlisting}[style=pystyle, caption={Yêu cầu tuyển dụng: FlashText tích hợp 8 taxonomy}, label={lst:clean-requirement}]
_TAXONOMY_TO_OUTPUT_COL = {
    "program_lang":       "require_programming_languages",
    "framework":          "require_frameworks",
    "tools":              "require_tools",
    "cloud_skill":        "require_cloud_skills",
    "knowledge":          "require_knowledge",
    "domain":             "require_domain_knowledge",
    "language":           "require_foreign_languages",
    "domain_university":  "require_domain_university",
}

def apply_requirement_cleaning(df, taxonomies, text_col="requirements",
                               cleaned_col="requirements_cleaned",
                               extra_noise_patterns=None):
    df = df.with_columns(clean_text(text_col, extra_noise_patterns).alias(cleaned_col))
    # Accent-stripped view used ONLY for keyword matching
    df = df.with_columns(
        pl.col(cleaned_col).map_elements(remove_vietnamese_accents, return_dtype=pl.String)
        .str.to_lowercase().alias("_requirements_norm"))
    new_columns = []
    for tax_key, out_col in _TAXONOMY_TO_OUTPUT_COL.items():
        taxonomy_df = taxonomies.get(tax_key)
        if taxonomy_df is None or taxonomy_df.is_empty():
            new_columns.append(pl.lit([], dtype=pl.List(pl.String)).alias(out_col))
            continue
        extractor = build_extractor(taxonomy_df)
        new_columns.append(
            pl.col("_requirements_norm")
            .map_elements(
                lambda text, ex=extractor: extract_unique_labels(text, ex),
                return_dtype=pl.List(pl.String), skip_nulls=False)
            .alias(out_col))
    return df.with_columns(new_columns).drop("_requirements_norm")
\end{lstlisting}

Khi taxonomy rỗng/thiếu, output column vẫn được tạo với list rống
\texttt{[]} để Silver schema ổn định cho downstream.

% ============================================================
\subsubsection{Experience level: \texttt{clean\_exp\_level()}}

Tương tự salary nhưng đơn vị là năm. Phân loại: ``không yêu cầu''
$\to$ min=0, max=0; ``dưới X'' $\to$ min=0, max=X;
``trên X'' $\to$ min=X, max=null; ``X tháng'' $\to$ chia 12:

\begin{lstlisting}[style=pystyle, caption={Experience level: parse sang min/max năm}, label={lst:clean-exp}]
def clean_exp_level(df, column_name="experience_level",
                    min_col="min_exp_level", max_col="max_exp_level"):
    """Parse experience_level into 2 numbers (unit: YEARS).
    1. NULL/no digit -> MIN=0, MAX=NULL
    2. 'No req' / 'no exp' / 'fresher' -> MIN=0, MAX=0
    3. 'Under X' / 'below X' -> MIN=0, MAX=X
    4. 'Over X' / 'X+' -> MIN=X, MAX=NULL
    5. 'X months' -> divide by 12
    6. 'X - Y years' -> MIN=X, MAX=Y"""
    df = remove_accents_col(df, column_name, "exp_no_accent")
    # ... (flag + extract first/second number + unit_div)
    unit_div = pl.when(is_month).then(pl.lit(12.0)).otherwise(pl.lit(1.0))
    min_expr = (pl.when(raw.is_null()).then(zero_float)
                .when(is_no_req).then(zero_float)
                .when(is_less).then(zero_float)
                .otherwise(first_number / unit_div))
    # ...
    return df.with_columns([min_expr.alias(min_col), max_expr.alias(max_col)])
\end{lstlisting}

% ============================================================
\subsubsection{Hình thức công việc (Job type): \texttt{clean\_job\_type()}}

Cấu trúc đơn giản nhất — chuẩn hóa về 3 giá trị chuẩn (canonical):

\begin{lstlisting}[style=pystyle, caption={Hình thức công việc: 3 giá trị chuẩn hóa}, label={lst:clean-jobtype}]
def clean_job_type(df, column_name="job_type"):
    """Normalize to: 'Full-time', 'Part-time', or 'Other'."""
    df = remove_accents_col(df, column_name, "job_type_no_accent")
    t = pl.col("job_type_no_accent").str.to_lowercase()
    df = df.with_columns(
        pl.when(t.str.contains(r"(part|ban thoi gian)")).then(pl.lit("Part-time"))
        .when(t.str.contains(r"(full|toan thoi gian)")).then(pl.lit("Full-time"))
        .otherwise(pl.lit("Other")).alias("job_type"))
    return df.drop("job_type_no_accent")
\end{lstlisting}

% ============================================================
\subsection{Seed taxonomy system}
\label{subsec:seed-taxonomy}

% ============================================================
\subsubsection{Đọc dữ liệu nguồn (seeds): Ưu tiên Google Sheets, dự phòng bằng CSV}

Hàm \texttt{read\_seeds()} thử tải dữ liệu từ Google Sheets trước (nếu cấu hình thông tin xác thực sẵn sàng),
nếu gặp lỗi thì tự động dùng tệp CSV cục bộ dự phòng:

\begin{lstlisting}[style=pystyle, caption={Tải dữ liệu seed: ưu tiên Google Sheets, dự phòng bằng CSV cục bộ}, label={lst:read-seeds}]
def read_seeds(file_name_csv: str) -> pl.DataFrame:
    """Load seed data from Google Sheets first, then local CSV fallback."""
    if google_sheets_config_is_available():
        try:
            return read_worksheet_as_polars(worksheet_title_for_csv(file_name_csv))
        except GoogleSheetsError as exc:
            logger.warning("Failed to read %s from Google Sheets. Falling back: %s",
                           file_name_csv, exc)
    return read_local_seed(file_name_csv)
\end{lstlisting}

13 tệp taxonomy trong thư mục \texttt{silver/seeds/}:

\begin{longtable}{lll}
\toprule
\textbf{Tệp CSV} & \textbf{Mục đích} & \textbf{Cột đầu ra (Output column)} \\
\midrule
\endhead
\texttt{industry\_taxonomy.csv} & Ánh xạ ngành nghề & \texttt{job\_industry\_clean} \\
\texttt{benefit\_taxonomy.csv} & Ánh xạ phúc lợi & \texttt{benefits\_categories\_vi} \\
\texttt{program\_lang\_taxonomy.csv} & Ngôn ngữ lập trình & \texttt{require\_programming\_languages} \\
\texttt{framework\_taxonomy.csv} & Framework tuyển dụng & \texttt{require\_frameworks} \\
\texttt{tools\_taxonomy.csv} & Công cụ phát triển (Tool) & \texttt{require\_tools} \\
\texttt{cloud\_skill\_taxonomy.csv} & Kỹ năng điện toán đám mây (Cloud) & \texttt{require\_cloud\_skills} \\
\texttt{knowledge\_taxonomy.csv} & Kiến thức lý thuyết tổng quát & \texttt{require\_knowledge} \\
\texttt{domain\_taxonomy.csv} & Lĩnh vực chuyên ngành tuyển dụng & \texttt{require\_domain\_knowledge} \\
\texttt{language\_taxonomy.csv} & Yêu cầu ngoại ngữ & \texttt{require\_foreign\_languages} \\
\texttt{domain\_university\_taxonomy.csv} & Nhóm chuyên ngành đại học & \texttt{require\_domain\_university} \\
\texttt{company\_mapping.csv} & Ánh xạ quy chuẩn tên doanh nghiệp & \texttt{company\_name\_canonical} \\
\texttt{location\_mapping.csv} & Ánh xạ chuẩn hóa địa lý địa điểm & \texttt{clean\_location} \\
\texttt{tinh\_thanh.csv} & Danh sách tỉnh thành Việt Nam & (Trợ lý chuẩn hóa vị trí địa lý) \\
\bottomrule
\end{longtable}

% ============================================================
\subsubsection{Bộ trích xuất lai FlashText + Regex (FlashText + Regex hybrid extractor)}

Hàm \texttt{build\_extractor()} biên dịch (compile) taxonomy thành \texttt{HybridKeywordProcessor}
— sử dụng thư viện FlashText cho các từ khóa thuần túy (tốc độ $O(n)$), và sử dụng regex cho các biểu thức mẫu phức tạp
(ví dụ từ khóa nhạy cảm như \texttt{\textbackslash bIBM\textbackslash b}):

\begin{lstlisting}[style=pystyle, caption={Bộ trích xuất lai FlashText và regex kết hợp}, label={lst:build-extractor}]
def build_extractor(taxonomy_df, label_col="canonical_en", keywords_col="keywords"):
    """Compile hybrid (FlashText + regex) extractor from taxonomy DataFrame."""
    kp = KeywordProcessor(case_sensitive=False)
    regex_rules = []
    for row in taxonomy_df.iter_rows(named=True):
        raw_keywords = row.get(keywords_col) or ""
        label = row.get(label_col)
        if not raw_keywords or not label:
            continue
        for alternative in raw_keywords.split("|"):
            literal, regex_src = _classify_alternative(alternative)
            if literal is not None:
                kp.add_keyword(literal, label)
            elif regex_src is not None:
                regex_rules.append((re.compile(regex_src, re.IGNORECASE), label))
    return HybridKeywordExtractor(kp=kp, regex_rules=regex_rules)
\end{lstlisting}

% ============================================================
\subsubsection{Làm sạch văn bản: \texttt{clean\_text()}}

Hàm tiện ích dùng chung cho các trường thông tin mô tả chi tiết công việc, yêu cầu, và phúc lợi:
loại bỏ thẻ HTML (HTML tags), thực thể ký tự HTML (HTML entities), ký hiệu đầu dòng (bullets), các thành phần nhiễu đặc trưng của từng trang web, và thu gọn khoảng trắng dư thừa:

\begin{lstlisting}[style=pystyle, caption={Làm sạch văn bản: bỏ thẻ HTML, ký hiệu đầu dòng, và các thành phần nhiễu}, label={lst:clean-text}]
def clean_text(col, extra_noise_patterns=None):
    """Vectorised: strip HTML, optional site-noise, bullets, collapse whitespace."""
    expr = pl.col(col).cast(pl.String)
    expr = expr.str.replace_all(_HTML_TAG_RE, " ")
    expr = expr.str.replace_all(_HTML_ENTITY_RE, " ")
    for pattern in extra_noise_patterns or []:
        expr = expr.str.replace_all(pattern, " ")
    expr = expr.str.replace_all(_BULLET_RE, " ")
    expr = expr.str.replace_all(r"\s+", " ")
    return expr.str.strip_chars()
\end{lstlisting}

% ============================================================
\subsection{Đường ống xử lý đặc thù cho từng trang web (Site-specific pipelines)}
\label{subsec:site-pipelines}

Mỗi trang web thu thập (site) có một hàm \texttt{clean\_<site>\_jobs()} riêng nhằm liên kết chuỗi hàm làm sạch dùng chung (common cleaners)
theo thứ tự yêu cầu. Các đặc điểm phân biệt chính bao gồm:

\begin{lstlisting}[style=pystyle, caption={TopCV: dãy cleaner đầy đủ}, label={lst:clean-topcv}]
def clean_topcv_jobs(df):
    """TopCV has all fields; industry separated by comma."""
    df = drop_unecessary_cols(df)
    industry_taxonomy = read_seeds("industry_taxonomy.csv")
    df = clean_job_url(df, column_name="job_url")
    df = main_clean_company(df, column_name="company_name")
    df = process_job_title_pipeline(df, title_col="job_title")
    df = clean_location(df, column_name="location")
    df = apply_industry_cleaning(df, industry_taxonomy, sep=",")
    df = clean_topcv_description(df)
    df = clean_salary(df, column_name="salary")
    df = clean_topcv_benefit(df)
    df = clean_topcv_requirement(df)
    df = clean_company_size(df, column_name="company_size")
    df = clean_job_type(df); df = clean_exp_level(df)
    df = clean_education_lv(df); df = clean_job_position(df)
    df = clean_job_deadline(df)
    return df
\end{lstlisting}

ITViec bỏ qua 5 bước làm sạch cuối (job\_type, exp\_level, education\_lv,
job\_position, job\_deadline) do các cột thông tin này không tồn tại trong tệp JSONL thô thu thập được.
VietnamWorks bỏ qua bước \texttt{clean\_company\_size} (không tồn tại trường này) nhưng
chạy đầy đủ cả 5 bước làm sạch cuối. Các hàm làm sạch đặc thù \texttt{clean\_<site>\_benefit()}
và \texttt{clean\_<site>\_requirement()} hoạt động như một hàm bọc (wrapper) gọi đến
\texttt{apply\_benefit\_cleaning()} / \texttt{apply\_requirement\_cleaning()}
với cấu hình \texttt{extra\_noise\_patterns} riêng của từng trang web.

% ============================================================
\subsection{Kiểm định chất lượng dữ liệu Bronze (Great Expectations)}
\label{subsec:validation}

Trước khi tải lên (upload) tầng Bronze, tệp JSONL tạm thời (temp JSONL) có thể chạy qua quy trình kiểm định chất lượng (validation)
thông qua thư viện Great Expectations (GX) nhằm đảm bảo sự chuẩn xác về mặt dữ liệu:

\begin{lstlisting}[style=pystyle, caption={GX validation: ExpectColumnValuesToNotBeNull}, label={lst:gx-validate}]
def validate() -> bool:
    jsonl_path = _latest_jsonl(TEMP_DIR, FILE_PREFIX)
    df = pd.read_json(jsonl_path, lines=True)
    context = gx.get_context(mode="ephemeral")
    # ... (add pandas data source + batch definition)
    suite = _build_suite(context, df.columns.tolist())
    result = batch.validate(suite)
    _print_summary(result)
    return bool(result.success)

def _build_suite(context, columns):
    suite = context.suites.add(gx.ExpectationSuite(name="topcv_suite"))
    non_brand_row_condition = (
        f'not job_url.str.contains("{BRAND_URL_TOKEN}", na=False)')
    for column in columns:
        if column in NULLABLE_WHEN_BRAND_URL:
            # Brand profile page has no company_size -> skip null check
            suite.add_expectation(
                gx.expectations.ExpectColumnValuesToNotBeNull(
                    column=column, row_condition=non_brand_row_condition,
                    condition_parser="pandas"))
        else:
            suite.add_expectation(
                gx.expectations.ExpectColumnValuesToNotBeNull(
                    column=column, mostly=MIN_NON_NULL_RATIO))
    return suite
\end{lstlisting}

Tham số \texttt{MIN\_NON\_NULL\_RATIO} quy định tỷ lệ chấp nhận dữ liệu trống (null) nhỏ nhất
(không yêu cầu khắt khe 100\% dòng phải đầy đủ thông tin). Các trang thuộc nhóm hồ sơ thương hiệu (Brand profile page / company profile)
được áp dụng cơ chế ngoại lệ đặc thù để bỏ qua kiểm định chất lượng đối với trường \texttt{company\_size}.


% ============================================================
\subsection{Các phép biến đổi SQL ở tầng Gold (Gold SQL transformations)}
\label{subsec:gold-sql}

Khi tệp Silver dạng Parquet đã được lưu trữ trên S3, tầng Gold (Gold layer) thực hiện truy vấn trực tiếp qua hàm \texttt{read\_parquet()}
của DuckDB và xây dựng sơ đồ hình sao (star schema) thông qua chuỗi liên kết gồm 7 câu lệnh SQL.

% ============================================================
\subsubsection{Đọc Silver parquet: \texttt{read\_silver\_parquet\_sql()}}

\begin{lstlisting}[style=sqlstyle, caption={DuckDB read\_parquet với hive partitioning}, label={lst:read-parquet-sql}]
-- SILVER_PARQUET_GLOB = 's3://bucket/source_site=*/entity_name=*/*/*/*/*.parquet'
read_parquet(
    's3://thethien-lakehouse-lite-silver/source_site=*/.../*.parquet',
    hive_partitioning = true,
    filename = true,
    union_by_name = true
)
\end{lstlisting}

Cấu hình \texttt{hive\_partitioning=true} cho phép DuckDB tự động nhận biết các thông tin khóa phân vùng như
\texttt{source\_site}, \texttt{year}, \texttt{month}, \texttt{day}
trực tiếp từ đường dẫn tệp (path). Tham số \texttt{filename=true} giúp lưu giữ tên tệp gốc để dùng làm căn cứ loại trùng
(tệp mới hơn sẽ có thuộc tính \texttt{filename} lớn hơn).

% ============================================================
\subsubsection{Staging dedup: \texttt{build\_staging\_sql()}}

\begin{lstlisting}[style=sqlstyle, caption={Staging: dedup theo unique\_url + gán job\_id}, label={lst:staging-sql}]
CREATE OR REPLACE TABLE _staging AS
WITH source_rows AS (
    SELECT *, filename
    FROM read_parquet('s3://...', hive_partitioning=true, filename=true)
),
ranked AS (
    SELECT *,
        row_number() OVER (
            PARTITION BY COALESCE(unique_url, job_url)
            ORDER BY TRY_CAST(year AS INT) DESC,
                     TRY_CAST(month AS INT) DESC,
                     TRY_CAST(day AS INT) DESC,
                     filename DESC
        ) AS dedup_rank
    FROM source_rows
),
deduped AS (
    SELECT * EXCLUDE (filename, dedup_rank)
    FROM ranked WHERE dedup_rank = 1
)
SELECT
    ROW_NUMBER() OVER (ORDER BY COALESCE(unique_url, job_url)) AS job_id,
    * EXCLUDE (unique_url)
FROM deduped
\end{lstlisting}

Nếu cùng một khóa \texttt{unique\_url} xuất hiện trong nhiều ngày (do thu thập lại nhiều lần - re-crawl), bảng trung gian (staging)
chỉ lưu giữ dòng mới nhất (\texttt{ORDER BY year DESC, month DESC, day DESC,
filename DESC}). Sử dụng hàm \texttt{TRY\_CAST} để chuyển đổi khóa phân vùng từ chuỗi (string)
sang số nguyên (INT) nhằm tránh bị gián đoạn lỗi nếu gặp các giá trị không hợp lệ.

% ============================================================
\subsubsection{Dedup key: \texttt{build\_unique\_url\_expression()}}

\begin{lstlisting}[style=sqlstyle, caption={Unique URL dedup key (SQL)}, label={lst:unique-url-sql}]
-- Strip query/fragment; ITViec strips numeric suffix (-12345.html)
CASE
    WHEN lower(trim(CAST(source_site AS VARCHAR))) = 'itviec'
        THEN regexp_replace(
            regexp_replace(CAST(job_url AS VARCHAR), '[?#].*$', ''),
            '-[0-9]+/?$', '')
    ELSE regexp_replace(CAST(job_url AS VARCHAR), '[?#].*$', '')
END AS unique_url
\end{lstlisting}

Nếu trong tầng Silver đã có sẵn cột \texttt{unique\_url} (được khởi tạo bởi hàm \texttt{clean\_job\_url()}),
hệ thống sẽ ưu tiên sử dụng cấu trúc \texttt{COALESCE(NULLIF(trim(unique\_url), ''), <derived>)} để nhận giá trị đã làm sạch đó.

% ============================================================
\subsubsection{Dimension date: \texttt{build\_dim\_date\_sql()}}

\begin{lstlisting}[style=sqlstyle, caption={dim\_date: generate\_series 2023--2026}, label={lst:dim-date-sql}]
CREATE OR REPLACE TABLE gold.dim_date AS
WITH calendar AS (
    SELECT CAST(d AS DATE) AS full_date
    FROM generate_series(DATE '2023-01-01', DATE '2026-05-01', INTERVAL 1 DAY) AS g(d)
)
SELECT
    CAST(strftime(full_date, '%Y%m%d') AS INTEGER) AS date_key,
    full_date, year(full_date) AS year, quarter(full_date) AS quarter,
    month(full_date) AS month, monthname(full_date) AS month_name,
    day(full_date) AS day, isodow(full_date) AS day_of_week,
    dayname(full_date) AS day_name, week(full_date) AS week_of_year,
    (isodow(full_date) >= 6) AS is_weekend
FROM calendar ORDER BY full_date
\end{lstlisting}

% ============================================================
\subsubsection{Fact table: \texttt{build\_jobs\_sql()}}

\begin{lstlisting}[style=sqlstyle, caption={gold.jobs: fact table + date\_key}, label={lst:jobs-sql}]
CREATE OR REPLACE TABLE gold.jobs AS
SELECT
    job_id,
    job_url, clean_job_title, clean_company_name,
    min_monthly_salary, max_monthly_salary, ...,
    CAST(year AS INT) * 10000 + CAST(month AS INT) * 100 + CAST(day AS INT) AS date_key
FROM _staging
\end{lstlisting}

\texttt{date\_key} là integer \texttt{YYYYMMDD} để join với \texttt{dim\_date}.

% ============================================================
\subsubsection{Các bảng con liên kết (Child tables): UNNEST và UNION ALL}

Các bảng lưu trữ danh mục ngành nghề (industries) và phúc lợi (benefits) được tạo lập trực tiếp bằng phép biến đổi \texttt{UNNEST}:

\begin{lstlisting}[style=sqlstyle, caption={job\_industries: UNNEST list}, label={lst:industries-sql}]
CREATE OR REPLACE TABLE gold.job_industries AS
SELECT job_id, UNNEST(job_industry_clean) AS industry
FROM _staging
WHERE job_industry_clean IS NOT NULL AND len(job_industry_clean) > 0
\end{lstlisting}

Bảng yêu cầu tuyển dụng (Requirements) được xây dựng phức tạp hơn — cả 8 trường dữ liệu danh sách chuỗi (\texttt{List[String]}) đều được gộp chung vào
cùng một bảng \texttt{job\_requirements}, phân biệt bằng giá trị cột \texttt{requirement\_type}. Bản đồ ánh xạ (mapping) từ tên trường dữ liệu sang thuộc tính phân loại (discriminator)
và bảng con tương ứng được định nghĩa trong \texttt{LIST\_FIELD\_TO\_CHILD}:

\begin{lstlisting}[style=sqlstyle, caption={job\_requirements: Dùng phép UNION ALL kết hợp các nhánh dữ liệu}, label={lst:requirements-sql}]
-- Each branch: SELECT job_id, 'discriminator' AS requirement_type,
--              UNNEST(field) AS value FROM _staging WHERE ...
SELECT job_id, 'programming_language' AS requirement_type,
       UNNEST(require_programming_languages) AS value
FROM _staging WHERE require_programming_languages IS NOT NULL
    AND len(require_programming_languages) > 0
UNION ALL
SELECT job_id, 'framework' AS requirement_type,
       UNNEST(require_frameworks) AS value
FROM _staging WHERE require_frameworks IS NOT NULL
    AND len(require_frameworks) > 0
UNION ALL
-- ... (tool, cloud_skill, knowledge, domain_knowledge,
--      foreign_language, domain_university, special_keyword)
\end{lstlisting}

Tổng quan star schema chỉ trong Hình~\ref{fig:gold-star-schema}.

\begin{figure}[htbp]
  \centering
  % Screenshot: module3_diagrams.html diagram 3 (Gold star schema)
  \includegraphics[width=\textwidth]{module3_hinh3_3.png}
  \caption{Sơ đồ hình sao ở tầng Gold: loại trùng ở staging $\to$ bảng sự kiện fact jobs +
  bảng chiều dim\_date + các bảng con (industries, benefits, requirements).
  Cấu hình \texttt{LIST\_FIELD\_TO\_CHILD} ánh xạ mỗi trường dạng danh sách (List field) sang bảng con tương ứng.}
  \label{fig:gold-star-schema}
\end{figure}

% ============================================================
\subsubsection{7 lệnh SQL: \texttt{build\_all\_gold\_sql()}}

\begin{lstlisting}[style=pystyle, caption={Danh sách 7 câu lệnh Gold SQL}, label={lst:build-all-gold}]
def build_all_gold_sql(silver_columns):
    return [
        ("Staging table",    build_staging_sql(silver_columns, staging_table)),
        ("Date dimension",   build_dim_date_sql(dim_date_table)),
        ("Fact table",       build_jobs_sql(staging_table, jobs_table)),
        ("Industries table", build_industries_sql(staging_table, industries_table)),
        ("Benefits table",   build_benefits_sql(staging_table, benefits_table)),
        ("Requirements",     build_requirements_sql(staging_table, requirements_table)),
        ("Drop staging",     f"DROP TABLE IF EXISTS {staging_table}"),
    ]
\end{lstlisting}

Bảng staging đóng vai trò là bảng tạm thời — sau khi các bảng sự kiện (fact) và bảng con (child tables) được ghi dữ liệu hoàn tất,
bảng staging sẽ bị loại bỏ (drop) để tránh lãng phí không gian lưu trữ.

% ============================================================
\subsection{Bảng điều khiển báo cáo (Business dashboards)}
\label{subsec:dashboards}

Hai giao diện bảng điều khiển (dashboard) HTML được kết xuất trực tiếp từ các nguồn dữ liệu trên S3 bao gồm:

\begin{itemize}
  \item \textbf{Bronze dashboard} (\texttt{bronze\_dashboard.py}) — tổng
    hợp thống kê số lượng bản ghi, dung lượng lưu trữ (storage size) theo trang web nguồn và theo ngày từ Bronze
    S3. Báo cáo cung cấp các dạng biểu đồ (chart) phân tích lưu trữ theo site, theo ngày và số dòng dữ liệu hàng ngày.
  \item \textbf{Silver dashboard} (\texttt{silver\_dashboard.py}) — tổng
    hợp phân tích dữ liệu Silver theo khoảng thời gian tùy chọn, bao gồm cả bảng thống kê các ngành nghề chưa khớp taxonomy (unmapped industry)
    (hỗ trợ đánh giá và bổ sung danh mục taxonomy), thông tin doanh nghiệp, và mẫu dữ liệu phúc lợi/yêu cầu tuyển dụng.
\end{itemize}

Cả hai chương trình đều tải trực tiếp dữ liệu từ S3, tự động tạo mã HTML và ghi ra các tệp cục bộ hoặc tải trực tiếp lên S3.
Các dashboard được lên lịch vận hành tự động thông qua các Airflow DAG \texttt{bronze\_dashboard} và \texttt{silver\_dashboard}.


% ============================================================
\subsection{Các lưu ý kỹ thuật quan trọng (Gotchas)}
\label{subsec:processing-gotchas}

\begin{longtable}{p{0.42\textwidth} p{0.50\textwidth}}
\toprule
\textbf{Vấn đề kỹ thuật (Gotcha)} & \textbf{Giải thích chi tiết} \\
\midrule
\endhead
Thứ tự các hàm làm sạch bắt buộc & Hàm \texttt{clean\_2} phải được gọi trước \texttt{clean\_1}
trong chuỗi làm sạch tên công ty; ngoài ra, \texttt{clean\_location} và \texttt{apply\_industry\_cleaning}
sẽ xóa các cột dữ liệu gốc để giải phóng bộ nhớ, nên không hàm làm sạch nào phía sau được tham chiếu đến cột gốc \texttt{location} hoặc
\texttt{job\_industry}. \\
Schema đồng bộ từ dataclass & Cấu trúc \texttt{SILVER\_SCHEMA} được tự động tạo từ lớp dữ liệu
\texttt{SilverJobItem} thông qua hàm \texttt{silver\_schema\_to\_polars()}.
Khi thêm hoặc bớt trường dữ liệu, chỉ cần cập nhật dataclass mà không cần chỉnh sửa chuỗi schema thủ công. \\
Bỏ dấu để so khớp từ khóa & Quy trình so khớp từ khóa (keyword matching) chạy trên định dạng viết thường + không dấu,
nhưng giá trị trích xuất thực tế vẫn giữ nguyên dấu nguyên bản của từ (ví dụ \texttt{benefits\_text\_clean}
giữ đúng chữ ``tiếng Việt''). \\
Taxonomy trống sinh danh sách rỗng & Khi một tệp taxonomy trống hoặc thiếu, cột dữ liệu đầu ra tương ứng
vẫn được sinh ra với giá trị mặc định là danh sách rỗng \texttt{[]} để bảo vệ tính nhất quán của Silver schema cho các bước xử lý sau. \\
ITViec bỏ qua 5 bước làm sạch & Các trường \texttt{job\_type}, \texttt{exp\_level},
\texttt{education\_lv}, \texttt{job\_position}, \texttt{job\_deadline}
không tồn tại trong tệp JSONL của ITViec. Việc cố tình gọi các hàm này sẽ gây lỗi \texttt{ColumnNotFoundError}. \\
VNWorks không chứa company\_size & Tệp JSONL thu thập từ VietnamWorks hoàn toàn không chứa trường dữ liệu
\texttt{company\_size}. \\
Chạy lại không ghi đè dữ liệu & Bộ đọc dữ liệu Silver ưu tiên chọn tệp có nhãn thời gian \texttt{LastModified}
mới nhất cho mỗi ngày. Việc chạy lại quy trình chỉ tải lên tệp mới mà không ghi đè hoặc xóa bỏ tệp cũ. \\
Bảng staging của Gold là bảng tạm & Bảng staging đóng vai trò trung gian và bị loại bỏ (DROP) hoàn toàn sau khi các bảng sự kiện (fact) và bảng con đã được dựng xong. Việc chạy lại quy trình Gold sẽ tạo lại bảng staging từ đầu (sử dụng câu lệnh CREATE OR REPLACE). \\
\texttt{TRY\_CAST} cho khóa phân vùng & Các trường Year/month/day lấy từ đường dẫn tệp S3 ở dạng chuỗi văn bản.
Sử dụng hàm \texttt{TRY\_CAST} giúp chuyển đổi an toàn sang kiểu số nguyên (INT) nhằm tránh bị gián đoạn lỗi nếu gặp các giá trị không hợp lệ. \\
Theo dõi các giá trị chưa khớp & Sử dụng trường \texttt{job\_industry\_unmapped} và bảng liên kết
\texttt{job\_industries} để ghi nhận các ngành nghề lạ, giúp nhóm phát triển đánh giá chất lượng và cập nhật taxonomy kịp thời. \\
Cờ \texttt{-{}-no\_save} và \texttt{-{}-export\_parquet} & Cờ \texttt{-{}-no\_save} giúp
chạy thử nghiệm cục bộ không lưu lên S3; cờ \texttt{-{}-export\_parquet} kết xuất tệp parquet cục bộ phục vụ gỡ lỗi trực quan. \\
GX validation là tùy chọn & Quy trình kiểm định chất lượng dữ liệu chạy trước khi tải lên Bronze và không bắt buộc.
Tham số \texttt{MIN\_NON\_NULL\_RATIO} giúp chấp nhận một tỷ lệ lỗi null nhỏ để tránh dừng tiến trình không cần thiết. \\
\bottomrule
\end{longtable}

% ============================================================
\subsection{Lệnh chạy}
\label{subsec:processing-commands}

\begin{lstlisting}[style=pystyle, caption={Các lệnh chạy và vận hành Processing Layer}, label={lst:processing-commands}]
# 1. Validation (optional, before Bronze upload)
python -m src.storage_layer.MinIO_S3.layer.local_temp.validation.topcv_validate
python -m src.storage_layer.MinIO_S3.layer.local_temp.validation.itviec_validate
python -m src.storage_layer.MinIO_S3.layer.local_temp.validation.vnworks_validate

# 2. Silver cleaning (per site, specify date range)
python -m src.storage_layer.MinIO_S3.layer.silver.cleaning.clean_topcv.main_process \
    --from_date 2025-01-01 --to_date 2025-01-31

# 2b. Silver dry run (no upload)
python -m src.storage_layer.MinIO_S3.layer.silver.cleaning.clean_itviec.main_process \
    --from_date 2025-01-01 --to_date 2025-01-31 --no_save

# 2c. Silver + debug parquet dump
python -m src.storage_layer.MinIO_S3.layer.silver.cleaning.clean_vietnamworks.main_process \
    --from_date 2025-01-01 --to_date 2025-01-31 --export_parquet

# 3. Gold (MotherDuck star schema)
python -m src.storage_layer.MotherDuck.scripts.load_silver_to_gold

# 4. Supabase serving
python -m src.storage_layer.Supabase.scripts.load_silver_to_supabase \
    --from_date 2025-01-01 --to_date 2025-01-31

# 5. Business dashboards
python -m src.monitoring_layer.business.bronze_dashboard
python -m src.monitoring_layer.business.silver_dashboard
\end{lstlisting}
